\scanpage{139}\mspage{13}
% archive sheet number: 138
muni des flèches
\[
i_\alpha^{C_0}:\mathrm{id}_{\uline{\mathrm{Fl}}\,M}\longrightarrow\unsure{P_\alpha^{C_0}|p_\alpha^{C_0}}
\]
dépend fonctoriellement de $C_0$, quand on considère $C_0$ comme un élément de l'ensemble ordonné des \uline{petites} parties de $\mathrm{Fl}\,M$, \unsure{formées|formée} de flèches \unsure{coquarrables} (donc de l'ens ordonné
\[
\unsure{\mathscr P|\mathcal P}_{\unsure{\mathcal U|U}\text{-petites}}(\mathrm{Fl}_{\mathrm{coquar.}}(M))
\]
)\unsure{;|,} Si on a une inclusion
\[
C_0\subset C_1\tag*{(2.1.)}
\]
on en déduit \inserted{pour \unsure{tt} ordinal $\alpha$,} \deleted{$X_\alpha^{C_0}\unsure{\to|\simeq}X_\alpha$} un hom. fonctoriel
\[
u_\alpha=u_{\unsure{\alpha,f|\alpha,0|\alpha}}^{C_1,C_0}:X_\alpha^{C_0}\mathrel{\text{---}}X_\alpha^{C_1},\tag*{(2.2.)}
\]
rendant commutatifs les carrés déduits des $f:X\to Y$ de $M$
\diagram{\begin{tikzcd}[ampersand replacement=\&,column sep=large,row sep=large,execute at end picture={\node[anchor=east] at ([xshift=-30pt]current bounding box.west) {(2.3.)};}]
X\arrow[r,dash,"i_\alpha^{C_0}(f)"]\arrow[d,"\mathrm{id}"']\&X_\alpha^{C_0}(f)\arrow[r,phantom,"\unsure{-|\to}" description,"p_\alpha^{C_0}(f)"]\arrow[d,"u_\alpha(f)"]\&Y\arrow[d,dash,"\mathrm{id}"]\\
X\arrow[r,dash,"i_\alpha^{C_1}(f)"]\&X_\alpha^{C_1}(f)\arrow[r,"p_\alpha^{C_1}(f)"]\&Y
\end{tikzcd}}\unsure{.}
et compatible avec les morphismes de transition pour $\beta\unsure{\geq|>}\alpha$, \unsure{i.e.|c.à.d.}
\diagram{\begin{tikzcd}[ampersand replacement=\&,column sep=huge,row sep=large,execute at end picture={\node[anchor=east] at ([xshift=-30pt]current bounding box.west) {(2.4.)};}]
X_\alpha^{C_0}\arrow[r,"\mathrm{trans}_{\beta,\alpha}^{C_0}"]\arrow[d,"u_\alpha"']\&X_\beta^{C_0}\arrow[d,dash,"u_\beta"]\\
X_\alpha^{C_1}\arrow[r,"\mathrm{trans}_{\beta,\alpha}^{C_1}"']\&X_\beta^{C_1}
\end{tikzcd}}\unsure{.}
